\section{本章小结}
\label{sec:weak-summary}

章首的两个例子，$\operatorname{Bin}(n,2/n)$趋于$\operatorname{Poisson}(2)$与格点均匀分布趋于$\operatorname{Unif}(0,1)$，现在都有了确切的含义：对每个有界连续函数，分布下的平均收敛。这个定义只比较分布，不要求随机变量在同一空间上，也不断言样本彼此接近；在共同空间上，依概率收敛推出依分布收敛，反之只在极限为常量时成立。有界保证读数存在，连续让质点位置的微小移动只引起读数的微小变化。

本章的工具按承担的任务可以分成几类，各有各的条件，见表~\ref{tab:weak-summary}。

\begin{BookInlineTable}
 \centering
 \caption{弱收敛的工具及其条件}
 \label{tab:weak-summary}
 \begin{tabular}{@{}lll@{}}
 \toprule
 任务 & 工具 & 关键条件\\
 \midrule
 事件概率 & Portmanteau定理 & 极限在事件边界上无质量\\
 实轴上的判据 & 分布函数判据 & 只在极限分布函数的连续点\\
 极限的存在 & 胎紧与弱收敛子列 & 远处质量一致地小\\
 极限的识别 & 特征函数唯一性 & 无\\
 收敛的判别 & 连续性定理 & 逐点极限在原点连续\\
 另一种表示 & 矩母函数 & 在$0$的邻域上有限\\
 映射的传递 & 连续映射定理 & 极限不在不连续点集上\\
 联合极限 & Cramér--Wold判据 & 检查所有方向的投影\\
 常量替换 & Slutsky定理 & 一个分量依概率趋于常量\\
 期望的极限 & 一致可积 & 远处质量乘位置一致地小\\
 \bottomrule
 \end{tabular}
\end{BookInlineTable}

几组容易混淆的区别值得再强调。胎紧控制远处有多少质量，一致可积还要乘上位置：前者决定极限分布是否存在，后者决定期望是否随之收敛。特征函数总存在，识别分布不需要任何条件；矩母函数只在零点邻域有限时才能识别，对数正态分布说明所有矩有限也不够。边缘收敛不推出联合收敛，坐标方向的投影也不能代替所有方向；只有一个分量的极限是常量时，配对问题才自动消失。

这些工具中，弱收敛的证明常常分两步完成：用胎紧或特征函数在原点的连续性保证质量不逃逸，再用特征函数的唯一性识别极限。下一章的中心极限定理正是这样证明的：把标准化和的特征函数写成乘积，用二阶展开和乘积之差的引理求出高斯极限，再由连续性定理得到依分布收敛；多元情形则用Cramér--Wold判据化为各个方向。之后各章比较分布的距离时，也直接沿用本章的弱收敛，不再重新建立。
